\documentclass[10pt,a4paper]{amsart}
\usepackage{amsmath,amssymb,amsthm,graphicx}
\usepackage{verbatim}
\newenvironment{reference}{\comment}{\endcomment}
%\newenvironment{reference}{}{}
\newenvironment{slogan}{\comment}{\endcomment}
\newenvironment{history}{\comment}{\endcomment}
% For commutative diagrams we use Xy-pic
\usepackage[all]{xy}

% We use 2cell for 2-commutative diagrams.
\xyoption{2cell}
% Theorem environments.
%
\theoremstyle{plain}
\newtheorem{theorem}[subsection]{Theorem}
\newtheorem{proposition}[subsection]{Proposition}
\newtheorem{lemma}[subsection]{Lemma}

\theoremstyle{definition}
\newtheorem{definition}[subsection]{Definition}
\newtheorem{example}[subsection]{Example}
\newtheorem{exercise}[subsection]{Exercise}
\newtheorem{situation}[subsection]{Situation}

\theoremstyle{remark}
\newtheorem{remark}[subsection]{Remark}
\newtheorem{remarks}[subsection]{Remarks}

\numberwithin{equation}{subsection}

% Macros
%
\def\lim{\mathop{\mathrm{lim}}\nolimits}
\def\colim{\mathop{\mathrm{colim}}\nolimits}
\def\Spec{\mathop{\mathrm{Spec}}}
\def\Hom{\mathop{\mathrm{Hom}}\nolimits}
\def\Ext{\mathop{\mathrm{Ext}}\nolimits}
\def\SheafHom{\mathop{\mathcal{H}\!\mathit{om}}\nolimits}
\def\SheafExt{\mathop{\mathcal{E}\!\mathit{xt}}\nolimits}
\def\Sch{\mathit{Sch}}
\def\Mor{\mathop{\mathrm{Mor}}\nolimits}
\def\Ob{\mathop{\mathrm{Ob}}\nolimits}
\def\Sh{\mathop{\mathit{Sh}}\nolimits}
\def\NL{\mathop{N\!L}\nolimits}
\def\CH{\mathop{\mathrm{CH}}\nolimits}
\def\proetale{{pro\text{-}\acute{e}tale}}
\def\etale{{\acute{e}tale}}
\def\QCoh{\mathit{QCoh}}
\def\Ker{\mathop{\mathrm{Ker}}}
\def\Im{\mathop{\mathrm{Im}}}
\def\Coker{\mathop{\mathrm{Coker}}}
\def\Coim{\mathop{\mathrm{Coim}}}

% Boxtimes
%
\DeclareMathSymbol{\boxtimes}{\mathbin}{AMSa}{"02}

%
% Macros for moduli stacks/spaces
%
\def\QCohstack{\mathcal{QC}\!\mathit{oh}}
\def\Cohstack{\mathcal{C}\!\mathit{oh}}
\def\Spacesstack{\mathcal{S}\!\mathit{paces}}
\def\Quotfunctor{\mathrm{Quot}}
\def\Hilbfunctor{\mathrm{Hilb}}
\def\Curvesstack{\mathcal{C}\!\mathit{urves}}
\def\Polarizedstack{\mathcal{P}\!\mathit{olarized}}
\def\Complexesstack{\mathcal{C}\!\mathit{omplexes}}
% \Pic is the operator that assigns to X its picard group, usage \Pic(X)
% \Picardstack_{X/B} denotes the Picard stack of X over B
% \Picardfunctor_{X/B} denotes the Picard functor of X over B
\def\Pic{\mathop{\mathrm{Pic}}\nolimits}
\def\Picardstack{\mathcal{P}\!\mathit{ic}}
\def\Picardfunctor{\mathrm{Pic}}
\def\Deformationcategory{\mathcal{D}\!\mathit{ef}}
\usepackage{cleveref}
\crefname{originallemma}{original lemma}{original lemmas}
\Crefname{originallemma}{Original lemma}{Original lemmas}
\makeatletter
\newlabel{lemma-equivalence-finite-length}{{lemma-equivalence-finite-length}{source}{}{sourcepointer-1}{}}
\newlabel{lemma-equivalence-finite-length@cref}{{[originallemma][0][]lemma-equivalence-finite-length}{[1][0][]source}{}{}{}}
\makeatother
\crefname{originalproposition}{original proposition}{original propositions}
\Crefname{originalproposition}{Original proposition}{Original propositions}
\makeatletter
\newlabel{proposition-matlis}{{proposition-matlis}{source}{}{sourcepointer-2}{}}
\newlabel{proposition-matlis@cref}{{[originalproposition][0][]proposition-matlis}{[1][0][]source}{}{}{}}
\makeatother
\title{The dualizing functor on finite length modules over a local ring is the linear dual into the injective hull}
\author{The Stacks Project Authors}
\date{stacks-project:dualizing.tex:local-snapshot-20261012; GNU FDL 1.2 or later}
\begin{document}
\maketitle
\noindent\textit{Editorial note.} All delivered mathematical text is the author's own, copied from the pinned Stacks Project chapter source sources/stacks-dualizing/pinned-original-source.tex; no statement or proof step was rewritten, reordered or removed.\par\medskip
\noindent\textit{Editorial note.} Preamble provenance: the Stacks Project keeps one single shared preamble for every chapter; this card takes it directly from the staged sources/stacks-dualizing/preamble.tex as three verbatim spans (lines 8-12, 14-18 and 154-217), so no derived file was written and nothing was added or redefined. The wrapper class and the base packages are supplied by the assembler.\par\medskip
\noindent\textit{Editorial note.} Bibliography provenance: neither the counted statement nor its proof contains a citation (this lane rescans the whole span before writing the card), so the card carries no bibliography and no entry was reconstructed.\par\medskip
\noindent\textit{Editorial note.} Selection provenance: this card was opened under the file-based control contract (CONTROL.md), with the mechanical pre-card gate run on its statement and proof interval before the card was written.\par\medskip
\noindent\textit{Allowed standard prerequisites (editorial).}\begin{itemize}
\item Matlis duality: the functor into the injective hull induces an anti-equivalence on finite length modules and is its own inverse. Invoked inside the counted proof (source lines 3220-3222 and 3227). Gives the anti-equivalence and the double-dual identity.
\item The statement that an autoequivalence of the category of finite length modules is isomorphic to the identity. Invoked inside the counted proof (source lines 3223-3225). Identifies the composed functor with the identity.
\end{itemize}\medskip
\ifdefined\hypertarget\hypertarget{sourcepointer-1}{}\fi
\noindent\textit{Source reference.} Original lemma lemma-equivalence-finite-length: sources/stacks-dualizing/pinned-original-source.tex, source lines 3173--3173. Navigation pointer only: the counted text refers to the Stacks item lemma-equivalence-finite-length, which is an existing labeled result of this same chapter and is not part of this card. The pointer carries the target's own label as its number and its definition line as its span; its content is not reproduced here.\par\medskip
\ifdefined\hypertarget\hypertarget{sourcepointer-2}{}\fi
\noindent\textit{Source reference.} Original proposition proposition-matlis: sources/stacks-dualizing/pinned-original-source.tex, source lines 1185--1185. Navigation pointer only: the counted text refers to the Stacks item proposition-matlis, which is an existing labeled result of this same chapter and is not part of this card. The pointer carries the target's own label as its number and its definition line as its span; its content is not reproduced here.\par\medskip
\begin{lemma}
\label{lemma-dualizing-finite-length}
Let $(A, \mathfrak m, \kappa)$ be a Noetherian local
ring with normalized dualizing complex $\omega_A^\bullet$.
Let $E$ be an injective hull of $\kappa$. Then there exists
a functorial isomorphism
$$
R\Hom_A(N, \omega_A^\bullet) = \Hom_A(N, E)[0]
$$
for $N$ running through the finite length $A$-modules.
\end{lemma}

\begin{proof}
By induction on the length of $N$ we see that $R\Hom_A(N, \omega_A^\bullet)$
is a module of finite length sitting in degree $0$. Thus
$R\Hom_A(-, \omega_A^\bullet)$ induces an anti-equivalence
on the category of finite length modules. Since the same is true
for $\Hom_A(-, E)$ by Proposition \ref{proposition-matlis} we see that
$$
N \longmapsto \Hom_A(R\Hom_A(N, \omega_A^\bullet), E)
$$
is an equivalence as in Lemma \ref{lemma-equivalence-finite-length}.
Hence it is isomorphic to the identity functor.
Since $\Hom_A(-, E)$ applied twice is the identity
(Proposition \ref{proposition-matlis}) we obtain
the statement of the lemma.
\end{proof}

\end{document}
